\documentclass{ieeeaccess}
\usepackage{cite}
\usepackage{amsmath,amssymb,amsfonts}
\usepackage{algorithmic}
\usepackage{graphicx}
\usepackage{textcomp}
\def\BibTeX{{\rm B\kern-.05em{\sc i\kern-.025em b}\kern-.08em
    T\kern-.1667em\lower.7ex\hbox{E}\kern-.125emX}}
\def\xfigwd{0pt}
\begin{document}
\history{Date of publication xxxx 00, 0000, date of current version xxxx 00, 0000.}
\doi{10.1109/ACCESS.2024.DOI}

\title{Innovation Capability and Startup Survival: An Empirical Study
of Early-Stage Ventures in India}

\author{\uppercase{Author One}\authorrefmark{1},
\uppercase{Author Two}\authorrefmark{1},
\uppercase{Author Three}\authorrefmark{1},
\uppercase{Author Four}\authorrefmark{1}}

\address[1]{Department of Industrial Engineering \& Management,
RV College of Engineering, Bengaluru -- 560059, India
(e-mail: \{author1, author2, author3, author4\}@rvce.edu.in)}

\tfootnote{This work was conducted as part of the VI Semester Experiential
Learning (EL) Project, Department of Industrial Engineering \& Management,
RV College of Engineering, Bengaluru, India.}

\markboth
{Author \headeretal: Innovation Capability and Startup Survival in India}
{Author \headeretal: Innovation Capability and Startup Survival in India}

\corresp{Corresponding author: Author One (e-mail: author1@rvce.edu.in).}

\begin{abstract}
Of every ten startups that launch in India today, roughly seven will cease
operations before their fifth anniversary---a failure rate that has
persisted despite expanding government support, institutional incubation
infrastructure, and record venture capital inflows \cite{b1,b2}. Innovation
capability, broadly defined as an organisation's capacity to commercialise
novel products, improve internal processes, and develop distinctive
market-entry strategies, is widely cited as a determinant of venture
survival; yet the literature has engaged this claim with considerable
conceptual looseness, treating innovation as a single undifferentiated
variable and thereby obscuring the distinct mechanisms by which product
innovation capability (PIC), process innovation capability (PROC), and
marketing innovation capability (MIC) each shape survival prospects.
No study, to date, has examined this disaggregated relationship in the
Indian context using a structural modelling approach. This paper proposes
a model---grounded in Dynamic Capabilities Theory \cite{b10} and the
Resource-Based View (RBV) \cite{b11}---in which PIC, PROC, and MIC
influence startup survival (SS) both directly and indirectly through
competitive advantage (CA) as a mediating variable, with founder
experience (FE) and funding stage (FS) specified as moderators that
condition the strength of these effects. Hypothesis testing will employ
Partial Least Squares Structural Equation Modelling (PLS-SEM) on data
collected from 200--300 DPIIT-registered Indian startup founders via a
structured five-point Likert-scale questionnaire; reliability thresholds
of Cronbach's $\alpha > 0.7$, composite reliability $> 0.7$, and average
variance extracted (AVE) $> 0.5$ will be applied throughout. The study
offers the first SEM-based, type-disaggregated empirical analysis of
innovation capability and startup survival in India, generating insights
relevant to founders, investors, and ecosystem policymakers alike.
\end{abstract}

\begin{keywords}
Competitive advantage, dynamic capabilities, entrepreneurship, India,
innovation capability, PLS-SEM, startup survival
\end{keywords}

\titlepgskip=-15pt

\maketitle

% ==========================================================
\section{Introduction}
\label{sec:introduction}
% ==========================================================

\PARstart{S}{omeone} founding a technology venture in Bengaluru,
Mumbai, or Hyderabad today enters what is nominally the world's
third-largest startup ecosystem \cite{b2}. The odds, however, are
not favourable. Roughly 70\% of Indian startups shut down within five
years of inception, and across high-growth tech sectors attrition
climbs higher still \cite{b1}. Capital is more available than at any
prior point in India's entrepreneurial history---2021 alone saw over
\$42 billion in venture funding flow into Indian startups---yet failure
rates have not declined proportionately. That contradiction warrants
explanation.

Innovation capability is the most frequently invoked explanatory
variable in startup survival research. The intuition is transparent:
ventures able to develop differentiated products, refine operational
processes, and reconfigure market strategies continuously are better
equipped to outlast rivals and environmental shocks \cite{b3}. Broad
empirical support for this claim exists \cite{b4, b5}. The difficulty
is that most studies operationalise innovation as a single composite
score---an approach that conceals more than it reveals. Product
innovation, process innovation, and marketing innovation are not
interchangeable activities. Each draws on different organisational
resources, carries a distinct risk profile, and interacts differently
with market conditions \cite{b6}. A startup that excels at new-product
development but lacks process discipline may exhaust its runway before
reaching market scale; one that over-invests in operational efficiency
prematurely can sacrifice the customer-facing agility on which early
survival depends. Disaggregating the innovation construct is, therefore,
not a methodological nicety---it changes the substantive conclusions.

A second limitation of existing work is geographic. The overwhelming
bulk of empirical studies on innovation and startup survival originates
in North America or Western Europe \cite{b7}. Institutional context
matters in non-trivial ways: access to formal credit, the density of
supplier and distribution networks, regulatory predictability, and
prevailing attitudes toward failure all shape how innovation capabilities
translate into survival outcomes. Findings from Silicon Valley or the
German Mittelstand transfer imperfectly to markets where founders
frequently lack debt-financing collateral and a single regulatory
clearance can consume months of organisational attention.

A third gap concerns mechanism. Establishing that innovation capability
correlates with survival does not, by itself, explain \textit{how} that
relationship operates. Competitive advantage---the capacity to deliver
superior value relative to available alternatives---has been proposed
as the intervening construct \cite{b8}, but the mediation hypothesis
remains empirically untested for early-stage Indian ventures. The
degree to which founder experience and available funding moderate the
innovation--survival relationship has similarly received scant systematic
attention in this setting \cite{b9}.

This study addresses all three limitations within a single integrated
framework grounded in Dynamic Capabilities Theory \cite{b10} and the
RBV \cite{b11}. Nine research hypotheses are developed and will be
tested via PLS-SEM on a survey of 200--300 Indian startup founders,
examining how PIC, PROC, and MIC independently influence startup
survival, whether competitive advantage mediates these effects, and how
founder experience and funding stage condition the relationships.
Section~\ref{sec:lit} reviews the theoretical and empirical foundations.
Section~\ref{sec:framework} develops the conceptual model and hypotheses.
Section~\ref{sec:method} details the survey design and analytical
strategy. Sections~\ref{sec:results} and~\ref{sec:conclusion} present
expected findings and their implications.

% ==========================================================
\section{Literature Review}
\label{sec:lit}
% ==========================================================

\subsection{Innovation and Startup Survival}

Schumpeter's argument \cite{b12} that capitalist development proceeds
through recurring waves of ``creative destruction'' placed innovation---
not efficiency or scale---at the centre of competitive dynamics. Seven
decades of subsequent research have broadly validated this position at the
firm level: ventures that invest in innovation survive longer and grow
faster than those that do not \cite{b4}. What that research has been
slower to clarify is which \textit{kind} of innovation matters, and under
what conditions.

The Oslo Manual's taxonomy distinguishes product, process, and marketing
innovation as analytically separable activities that a firm may pursue
independently or in combination \cite{b3}. This distinction carries real
weight in the startup context. Product innovation requires sustained
R\&D effort and iterative customer-facing experimentation; process
innovation demands inward-facing operational discipline; marketing
innovation calls for go-to-market creativity and customer data
infrastructure. For a founding team with limited headcount and finite
capital, trade-offs among these activities are genuine and consequential.
Yet most empirical treatments of innovation capability in the survival
literature either collapse these dimensions into a composite index or
focus exclusively on product innovation, leaving process and marketing
effects largely unmeasured \cite{b6}.

\subsection{Dynamic Capabilities and Startup Performance}

The Dynamic Capabilities framework, introduced by Teece, Pisano, and
Shuen \cite{b10}, was developed to explain how firms sustain advantage
when competitive landscapes shift faster than static resource stocks
can respond. Three micro-level processes are central to the framework:
\textit{sensing}---scanning the environment for emerging threats and
opportunities; \textit{seizing}---reconfiguring resources to capture
identified opportunities; and \textit{reconfiguring}---transforming
existing competencies to maintain strategic alignment as conditions
change. For startups, all three processes are in near-constant operation:
the product roadmap pivots, the revenue model iterates, the founding
team restructures. Dynamic capabilities provide the theoretical vocabulary
for understanding why some ventures navigate these transitions successfully
and others do not \cite{b13}.

Eisenhardt and Martin \cite{b14} offered an important qualification: in
highly dynamic markets, these processes tend to converge toward
identifiable best practices rather than remaining wholly firm-specific
and idiosyncratic. Their reformulation is useful here. PIC and PROC can
be understood as observable, benchmarkable dynamic capabilities---ones
measurable through survey instruments and comparable across a startup
cohort---precisely as this study requires.

\subsection{Innovation Ecosystems and Social Capital}

No venture innovates in isolation, and the ecosystem within which a
startup is embedded shapes whether internal capabilities translate into
market value. Social capital---Putnam's \cite{b15} term for the aggregate
of network relationships, shared norms, and interpersonal trust that
enable coordinated action---has a dual relevance to startup innovation.
Structurally, dense networks supply referrals, warm introductions, and
early-adopter access. Cognitively, they accelerate the diffusion of tacit
knowledge that formal training and documentation cannot easily convey.

India's startup ecosystem has undergone substantial change since 2016.
The Startup India programme created a formal DPIIT-recognition
infrastructure; Tier-1 incubators such as NSRCEL Bengaluru, T-Hub
Hyderabad, and CIIE Ahmedabad have professionalised early-stage support;
and a maturing limited-partner base has extended venture capital
deeper into seed-stage rounds that were previously inaccessible to most
Indian founders \cite{b2}. Whether these improvements amplify the
impact of firm-level innovation capabilities on survival---or simply
substitute for weaker internal capabilities---is an open empirical
question that informs the moderating hypotheses in
Section~\ref{sec:framework}.

\subsection{Resource-Based View and Competitive Advantage}

Barney's \cite{b11} Resource-Based View rests on a deceptively simple
claim: performance differences between firms in the same industry
persist because strategically valuable resources are heterogeneously
distributed and imperfectly mobile. For a resource to generate sustained
competitive advantage, it must be valuable, rare, inimitable, and
non-substitutable---the VRIN criteria that have become standard
analytical vocabulary in strategic management \cite{b16}.

Innovation capability satisfies these criteria when it is embedded in
organisational routines and tacit knowledge rather than residing in
individual contributors or replicable procedures. A startup whose
capacity to innovate is inseparable from its decision-making culture,
iterative norms, and accumulated customer insight is considerably harder
to imitate than one whose advantage depends on a single algorithm or
registered patent. Alvarez and Busenitz \cite{b17} drew the
RBV--entrepreneurship connection explicitly, arguing that the
discovery and exploitation of opportunities---the defining act of
entrepreneurship---itself constitutes a VRIN resource. For the
survival question, the implication is that startups whose innovation
capabilities constitute genuine VRIN resources will systematically
outperform those whose advantage can be more readily matched.

\subsection{Competitive Advantage as a Mediating Mechanism}

Neither Dynamic Capabilities Theory nor the RBV specifies a direct
path from capability to survival---both frameworks operate through
competitive advantage as the proximate outcome. The logic is: capabilities
generate advantage; advantage generates performance; performance enables
survival. Empirical studies frequently short-circuit this chain, regressing
survival directly on capability measures and bypassing the mechanism
\cite{b4,b5}. The mediating link consequently remains largely unvalidated.

Camisón and Villar-López \cite{b18} provide partial evidence in a sample
of Spanish manufacturing firms, where innovation capability influenced
performance indirectly through competitive position. Their study is
important but limited in two respects: the sample consists of established
firms for whom competitive advantage is relatively stable; and the context
is European manufacturing, where market structures differ substantially
from those facing Indian tech startups still searching for product-market
fit. Whether the mediated pathway holds when competitive advantage is
transient and contested is the specific question this study pursues.

\subsection{Founder Characteristics as Boundary Conditions}

Identical innovation capabilities do not produce identical outcomes
across all startups---context modulates effect. Two of the most
theoretically grounded moderators are founder experience and funding stage.

On the experience side, Colombo and Grilli \cite{b19} showed in a
longitudinal study of Italian new-technology-based firms that founders
with prior industry-specific human capital grew their ventures
significantly faster---not because experienced founders innovated more,
but because they were more effective at converting innovation into
marketable output. The pathway runs through judgment: experienced founders
better understand customer purchasing behaviour, competitive positioning
dynamics, and the timing of market entry. In the Indian tech sector,
where investor due diligence consistently elevates founder quality as
the primary selection criterion, this moderation effect deserves
empirical scrutiny.

Funding stage operates through a different channel. External capital
relaxes the resource constraints that ordinarily attenuate the
innovation--survival relationship for early-stage ventures \cite{b20}.
A seed-funded startup can hire for execution; a pre-seed team must rely
on founder bandwidth alone. Hellmann and Puri \cite{b20} documented
that venture capital involvement accelerated operational
professionalisation in funded startups---a finding consistent with the
expectation that better-capitalised ventures extract greater survival
benefit from their innovation capabilities than their bootstrapped peers.

\subsection{Research Gaps and Study Positioning}

Five inter-related gaps emerge from this review. Indian startups have
not been the subject of a dedicated quantitative study linking innovation
capability to survival. The disaggregated effects of PIC, PROC, and MIC
on survival outcomes remain unmeasured. The mediating role of competitive
advantage has not been tested for early-stage venture populations. Founder
experience and funding stage have not been evaluated as moderators within
a structural model for this setting. And no PLS-SEM-based framework
integrating these constructs has been developed for the Indian startup
context. Each gap is directly addressed by the model described in
Section~\ref{sec:framework}.

% ==========================================================
\section{Conceptual Framework and Hypotheses}
\label{sec:framework}
% ==========================================================

\subsection{Conceptual Framework}

The proposed framework, illustrated in Fig.~\ref{fig:framework},
integrates Dynamic Capabilities Theory and the RBV into a single
explanatory structure. Three innovation capability constructs---PIC,
PROC, and MIC---serve as independent variables. Competitive advantage
(CA) is positioned as a mediator through which capabilities translate
into startup survival (SS). Founder experience (FE) and funding stage
(FS) function as moderators conditioning the strength of both the
capability--competitive advantage and capability--survival paths. Path
directions and expected signs are derived deductively from the
theoretical arguments developed in Section~\ref{sec:lit}.

\begin{figure}[htbp]
  \centering
  \fbox{\parbox{0.85\columnwidth}{\centering\vspace{1.5cm}
    [Conceptual Framework Diagram Placeholder]\\[0.2cm]
    PIC, PROC, MIC $\rightarrow$ CA $\rightarrow$ SS\\
    Moderators: FE, FS
    \vspace{1.5cm}}}
  \caption{Proposed conceptual framework linking disaggregated innovation
  capabilities, competitive advantage, and startup survival.}
  \label{fig:framework}
\end{figure}

\subsection{Hypothesis Development}

\textit{Direct effects on startup survival:}
Product innovation capability enables startups to offer differentiated
value propositions that address unmet customer needs, directly improving
market viability and revenue continuity \cite{b4}. Process innovation
capability reduces operational waste and strengthens the cost discipline
that sustains ventures through periods of low revenue \cite{b6}. Marketing
innovation capability opens novel customer acquisition and retention
channels that are difficult for resource-constrained incumbents to replicate
quickly \cite{b3}. Three hypotheses follow:

\begin{description}
  \item[H1:] Product innovation capability (PIC) has a significant
  positive effect on startup survival (SS).
  \item[H2:] Process innovation capability (PROC) has a significant
  positive effect on startup survival (SS).
  \item[H3:] Marketing innovation capability (MIC) has a significant
  positive effect on startup survival (SS).
\end{description}

\textit{Mediation by competitive advantage:}
The RBV's value-creation--value-capture logic implies that innovation
capabilities survive the competitive marketplace by first generating
a superior position relative to rivals; it is that position---not the
capability itself---that directly sustains the venture \cite{b8,b11}.
Modelling competitive advantage as a mediator makes this chain explicit:

\begin{description}
  \item[H4:] Competitive advantage (CA) mediates the relationship
  between innovation capabilities (PIC, PROC, MIC) and startup survival (SS).
\end{description}

\textit{Moderation by founder experience and funding stage:}
Human capital theory predicts that experienced founders extract more
survival benefit from a given level of innovation capability \cite{b9},
and resource-dependence arguments suggest the same for better-funded
ventures \cite{b20}:

\begin{description}
  \item[H5:] Founder experience (FE) positively moderates the
  relationship between innovation capabilities and startup survival.
  \item[H6:] Funding stage (FS) positively moderates the relationship
  between innovation capabilities and startup survival.
\end{description}

\textit{Innovation capabilities and competitive advantage:}
Competitive advantage can act as a mediator only if it is itself
a function of innovation capabilities. Establishing the
capability--advantage relationship is therefore a precondition for
the mediation test:

\begin{description}
  \item[H7:] Product innovation capability (PIC) has a significant
  positive effect on competitive advantage (CA).
  \item[H8:] Process innovation capability (PROC) has a significant
  positive effect on competitive advantage (CA).
  \item[H9:] Marketing innovation capability (MIC) has a significant
  positive effect on competitive advantage (CA).
\end{description}

% ==========================================================
\section{Research Methodology}
\label{sec:method}
% ==========================================================

\subsection{Research Design}

The study takes a deductive stance, deriving testable hypotheses from
established theory and subjecting them to empirical scrutiny through
primary survey data. A quantitative, cross-sectional design is adopted
in preference to a longitudinal approach---a trade-off that sacrifices
causal precision for practical feasibility, given the difficulty of
tracking startup cohorts over time and the near-certainty of differential
attrition among failing ventures \cite{b21}. The approach follows the
positivist tradition of quantitative entrepreneurship research, in which
constructs are operationalised through multi-item reflective scales and
structural relationships are assessed through second-generation
multivariate analysis.

\subsection{Target Population and Sampling}

Three inclusion criteria define the target population. Respondents must
hold a founding or co-founding role in the venture. The startup must be
formally recognised by the Department for Promotion of Industry and
Internal Trade (DPIIT) under the Startup India scheme. The venture must
have been incorporated no more than seven years before the survey date---
a ceiling consistent with DPIIT's own statutory definition of a startup
and with the early-growth phase during which the innovation--survival
relationship is most consequential. Sectoral scope is restricted to
technology-intensive industries (information technology, fintech,
healthtech, edtech, and agritech) where innovation capability variation
is sufficiently wide to permit structural model estimation.

Respondents will be reached through startup incubators, institutional
entrepreneurship cells, accelerator alumni networks, and LinkedIn
outreach. Purposive sampling is appropriate given the specificity of
the inclusion criteria; snowball referrals will supplement direct
outreach where response rates prove insufficient. Hair et al.'s
\cite{b22} ten-times rule implies a minimum of approximately 50 usable
responses for the present model structure; the target of 200--300
provides a substantial buffer and enables subgroup analyses across
funding stage categories.

\subsection{Measurement Instrument}

The questionnaire comprises five parts. Part~A collects demographic
data on the respondent and the venture: founding year, sector, current
funding stage, prior entrepreneurial and industry experience, and team
size. Parts~B through~D measure the three innovation capability constructs
using five items each, adapted from Romijn and Albaladejo \cite{b6} and
Gunday et al.\ \cite{b3} with language tailored to early-stage venture
settings. Part~E covers competitive advantage (four items, adapted from
Camisón and Villar-López \cite{b18}) and startup survival (five items
operationalised as perceived revenue continuity, customer retention
trajectory, and operational longevity relative to sector peers); founder
experience is measured using three items, while funding stage is treated
as a nominal moderator with five ordinal levels: bootstrapped,
angel-funded, seed-funded, Series~A or above, and grant-supported.

All multi-item measures use a five-point Likert scale anchored at
$1 =$ ``Strongly Disagree'' and $5 =$ ``Strongly Agree.'' A pilot
study involving 20 respondents will be conducted before main data
collection to identify problematic items and confirm scale properties.

\subsection{Analytical Strategy}

PLS-SEM via SmartPLS 4.0 \cite{b22} is selected over covariance-based
SEM for three reasons specific to this study: the model is
exploratory-predictive in orientation rather than purely confirmatory;
several constructs are expected to exhibit non-normal distributions
given the skewed nature of startup performance outcomes; and the
anticipated sample size, while adequate, falls short of the 400-plus
observations typically recommended for robust CB-SEM estimation.

Assessment follows the two-stage sequence recommended by Hair et al.\
\cite{b22}. The \textit{measurement model} is evaluated first: item
loadings must exceed 0.70; Cronbach's $\alpha$ and composite reliability
(CR) must exceed 0.70; and AVE must exceed 0.50 to establish convergent
validity. Discriminant validity is assessed via the
Heterotrait--Monotrait ratio (HTMT $< 0.85$), a more stringent test
than the traditional Fornell--Larcker criterion.

\textit{Structural model} evaluation proceeds through PLS bootstrapping
with 5,000 resamplings, generating $t$-statistics and 95\% confidence
intervals for each path coefficient (H1--H3, H7--H9). The mediation
test for H4 applies the bias-corrected bootstrap confidence interval
to the indirect effect: significance of the indirect effect alongside
a non-significant direct effect constitutes full mediation. Moderation
tests for H5 and H6 use the product-indicator interaction term method
within PLS-SEM; effect sizes are reported as Cohen's $f^2$ (small:
0.02, medium: 0.15, large: 0.35). Model fit is evaluated using the
Standardised Root Mean Square Residual (SRMR $< 0.08$).

Common method variance will be screened using Harman's single-factor
test; where the single-factor solution accounts for more than 50\% of
total variance, the common latent factor technique will be applied as
a post-hoc correction \cite{b21}.

% ==========================================================
% Remaining sections — to be completed
% ==========================================================

\section{Expected Results and Discussion}
\label{sec:results}

\subsection{Expected Findings per Hypothesis}

% TODO: H1--H3 direct effects discussion (~150 words)

\subsection{Mediation and Moderation Effects}

% TODO: H4--H6 mediation/moderation discussion (~100 words)

\subsection{Theoretical Implications}

% TODO: Link to Dynamic Capabilities Theory and RBV (~80 words)

\subsection{Practical Implications for the Indian Startup Ecosystem}

% TODO: Policy/founder/investor takeaways (~80 words)

\section{Conclusion}
\label{sec:conclusion}

\subsection{Summary of Objectives and Approach}

% TODO: Brief recap of research questions and method (~80 words)

\subsection{Theoretical and Practical Contributions}

% TODO: Key contributions to literature and practice (~70 words)

\subsection{Limitations and Future Research}

% TODO: Sample, cross-sectional design limitations; future longitudinal/qualitative directions (~60 words)

\appendices

\section*{Acknowledgment}

The authors thank the faculty of the Department of Industrial Engineering
\& Management, RV College of Engineering, for their guidance throughout
the Experiential Learning project.

\begin{thebibliography}{00}

\bibitem{b1} CB Insights, ``The top 12 reasons startups fail,''
CB Insights Research, New York, NY, USA, Tech. Rep., 2023.

\bibitem{b2} Department for Promotion of Industry and Internal Trade
(DPIIT), ``Annual report on Startup India,'' Ministry of Commerce and
Industry, Government of India, New Delhi, India, 2023.

\bibitem{b3} R. Gunday, G. Ulusoy, K. Kilic, and L. Alpkan, ``Effects of
innovation types on firm performance,'' \emph{Int. J. Prod. Econ.},
vol.~133, no.~2, pp.~662--676, Oct. 2011.

\bibitem{b4} A. Coad and R. Rao, ``Innovation and firm growth in
high-tech sectors: A quantile regression approach,'' \emph{Res. Policy},
vol.~37, no.~4, pp.~633--648, May 2008.

\bibitem{b5} M. Freel and R. Robson, ``Small firm innovation, growth and
performance: Evidence from Scotland and Northern England,''
\emph{Int. Small Bus. J.}, vol.~22, no.~6, pp.~561--575, 2004.

\bibitem{b6} H. A. Romijn and M. Albaladejo, ``Determinants of innovation
capability in small electronics and software firms in southeast England,''
\emph{Res. Policy}, vol.~31, no.~7, pp.~1053--1067, 2002.

\bibitem{b7} R. La Porta, F. Lopez-de-Silanes, A. Shleifer, and
R. Vishny, ``Legal determinants of external finance,''
\emph{J. Finance}, vol.~52, no.~3, pp.~1131--1150, 1997.

\bibitem{b8} M. E. Porter, \emph{Competitive Advantage: Creating and
Sustaining Superior Performance.} New York, NY, USA: Free Press, 1985.

\bibitem{b9} T. Åstebro and I. Bernhardt, ``The winner's curse of human
capital,'' \emph{Small Bus. Econ.}, vol.~24, no.~1, pp.~63--78, 2005.

\bibitem{b10} D. J. Teece, G. Pisano, and A. Shuen, ``Dynamic capabilities
and strategic management,'' \emph{Strategic Manage. J.}, vol.~18, no.~7,
pp.~509--533, Aug. 1997.

\bibitem{b11} J. B. Barney, ``Firm resources and sustained competitive
advantage,'' \emph{J. Manage.}, vol.~17, no.~1, pp.~99--120, 1991.

\bibitem{b12} J. A. Schumpeter, \emph{Capitalism, Socialism and Democracy.}
New York, NY, USA: Harper \& Brothers, 1942.

\bibitem{b13} D. J. Teece, ``Explicating dynamic capabilities: The nature
and microfoundations of (sustainable) enterprise performance,''
\emph{Strategic Manage. J.}, vol.~28, no.~13, pp.~1319--1350, 2007.

\bibitem{b14} K. M. Eisenhardt and J. A. Martin, ``Dynamic capabilities:
What are they?'' \emph{Strategic Manage. J.}, vol.~21, no.~10--11,
pp.~1105--1121, 2000.

\bibitem{b15} R. D. Putnam, \emph{Bowling Alone: The Collapse and Revival
of American Community.} New York, NY, USA: Simon \& Schuster, 2000.

\bibitem{b16} J. B. Barney, ``Looking inside for competitive advantage,''
\emph{Acad. Manage. Perspect.}, vol.~9, no.~4, pp.~49--61, 1995.

\bibitem{b17} S. A. Alvarez and L. W. Busenitz, ``The entrepreneurship of
resource-based theory,'' \emph{J. Manage.}, vol.~27, no.~6,
pp.~755--775, 2001.

\bibitem{b18} C. Camisón and A. Villar-López, ``Organizational innovation
as an enabler of technological innovation capabilities and firm
performance,'' \emph{J. Bus. Res.}, vol.~67, no.~1, pp.~2891--2902,
2014.

\bibitem{b19} M. G. Colombo and L. Grilli, ``Founders' human capital and
the growth of new technology-based firms: A competence-based view,''
\emph{Res. Policy}, vol.~34, no.~6, pp.~795--816, 2005.

\bibitem{b20} T. Hellmann and M. Puri, ``Venture capital and the
professionalisation of start-up firms: Empirical evidence,''
\emph{J. Finance}, vol.~57, no.~1, pp.~169--197, 2002.

\bibitem{b21} P. M. Podsakoff, S. B. MacKenzie, J.-Y. Lee, and
N. P. Podsakoff, ``Common method biases in behavioral research: A
critical review of the literature and recommended remedies,''
\emph{J. Appl. Psychol.}, vol.~88, no.~5, pp.~879--903, 2003.

\bibitem{b22} J. F. Hair, G. T. M. Hult, C. M. Ringle, and
M. Sarstedt, \emph{A Primer on Partial Least Squares Structural Equation
Modeling (PLS-SEM),} 3rd~ed. Thousand Oaks, CA, USA: SAGE, 2022.

\end{thebibliography}

\EOD

\end{document}
